\section{Tasks in AI Software Engineering}

\subsection{Code Generation}

Code generation is the task of generating code from a specification. In \textit{code completion} (e.g. Github Copilot), the specification takes the form of a pre-existing code snippet, and the goal is to complete the snippet. While code completion is perhaps the most successful paradigm, another popular one is \textit{natural language to code}, where the specification is a natural language description containing information such as the task description, input-output examples, libraries to use, and other requirements about the code. NL to code has seen success in specialized domains such as text to SQL.\ks{again we should using "logical complexity"}

In the industry, there have been many AI coding assistants such as \href{https://docs.cursor.com/composer}{Cursor Composer} and \href{https://codeium.com/windsurf}{Codeium's Windsurf Editor} that aim to serve as AI-driven IDEs. These systems blur the lines between the two paradigms. With the ultimate goal of understanding the user's intent and decreasing the burden of human programmers, these AI systems both automatically complete code when the intent is clear and use chat interfaces for the user to specify desired functionality.

Code generation can vary highly in difficulty, depending on scope and logical complexity. Today, large codebases still prove to be a challenge for state-of-the-art AI systems. \todo{elaborate}

\ks{Add contents on secure code generation.}
\todo{
\\ Constrained Decoding for Secure Code Generation
https://arxiv.org/abs/2405.00218
\\ PurpleLlama
https://github.com/meta-llama/PurpleLlama}

\ks{This section is a bit thin on related work.  It should be expanded to include more related work.}



\subsection{Code Transformation}

\textbf{Code Refactoring:} The task of code refactoring is to take a working implementation of a piece of software and rewrite parts of it while maintaining correctness. One challenge with this task is that success extends beyond functional correctness or metrics. While code refactoring often has a low logical complexity, it can be laborious in practice due to scope, as seemingly small refactors can propagate across the entire codebase. The goal is often to improve code maintainability, readability, or extensibility—qualities that can be inherently difficult to quantify and highly context-dependent. 

For instance, extracting shared functionality into helper methods presents trade-offs between modularity and cognitive complexity \citep{parnas1972criteria}. While there are no hard rules for when to extract functionality, one heuristic adopted by software engineers is the rule of three (``three strikes and you refactor'')--abstractions should only be used when a piece of code has been duplicated thrice. Because it can often be unclear what level of abstraction refactorings should be done at, completing a refactoring at a high autonomy level is also difficult. These challenges are further compounded by the need to understand implicit trade-offs customized to specific codebases, respect conventions, and reason about the long-term maintenance implications of structural changes. One benchmark for measuring this capability is RefactorBench \citep{gautam2024refactorbench}.

\textbf{Code Migration:} An incredibly resource-intensive (time and manual effort) task frequently affecting companies is transforming code to accommodate new language versions, framework updates, changing APIs and data types. A recent Scala 2.13 to 3 migration~\citep{scalamigrate} illustrates these challenges, documenting a year-long effort. Critical issues included the loss of macro annotations, broken type projections, incompatible libraries, and compiler performance degradation—all requiring innovative workarounds and architectural changes. Another example is COBOL, which powers 80\% of in-person financial services transactions and 95\% of ATM swipes while processing \$3 trillion in commerce a day \citep{Taulli_2020}. Although there are over 220 billion lines of COBOL code, there are less and less COBOL programmers, leading to many attempts to slowly migrate out of it \citep{sneed2001extracting, sellink2002restructuring, sneed2010migrating}.

Such high-value migrations are opportunities for AI-assisted automation to reduce cost and manual effort. Works at Google~\citep{nikolov2025google} and Amazon~\citep{omidvar2024evaluating} have explored LLM-driven solutions for simple but vast migrations.
% 
Google's system identifies locations for changes, generates edits with LLMs fine-tuned on internal code, and automatically validates changes through compilation and test execution.
% 
For a 500+M LoC project migrating IDs from 32-bit to 64-bit integers, their system produced 80\% AI-authored modifications and reduced engineering time by approximately 50\%. 

However, these approaches depend critically on the scope (number of files and systems affected) and logical complexity (semantic depth of required transformations). Current solutions may excel at migrations with high scope but modest logical demands (API migrations, type conversions) but struggle with changes across component boundaries~\citep{nikolov2025google}. Additionally, AI code migrations may require substantial human oversight due to their safety-critical and cross-system nature, motivating more work on Human-AI collaboration.

% \todo{At least 3 PL people I talked to mentioned the COBOL case study, probably worth to include it ;) \citep{sneed2001extracting, sellink2002restructuring, sneed2010migrating}}
% \todo{add code migration, deprecation, version upgrades}
% \ks{May be added another item on Code upgrade from one version to another.}

% Notes:
% talk about expensive and difficult, hence high-value for automation that AI brings
% 1. Some Py2 -> 3 migration examples
% 2. Scala 2.13 -> 3 migration pain points: https://blog.pierre-ricadat.com/scala-3-migration-report-from-the-field -- nice one!
% 3. Browser JS migrations?
% LLM for migration at google: https://research.google/blog/accelerating-code-migrations-with-ai/

\textbf{Code Translation (Transpilation):} Beyond migrating between versions, code translation transforms code from a source to a target language. In industry, this task can be motivated by several reasons, such as security and scalability concerns in legacy languages and avoiding the technical debt a project has accumulated over the years.
% 
Twitter\footnote{\url{https://www.infoq.com/news/2012/11/twitter-ruby-to-java/}} built its initial platform using Ruby on Rails, facilitating rapid development. However, as the user base expanded, performance and scalability issues arose. They migrated key components to Java and Scala, achieving a 3X latency drop. This transition required re-architecting the system to adapt Ruby's dynamic features to the statically typed environments of Java and Scala, exemplifying the complexities of large-scale code translation.

More recently, there has been a push to use translation as a proactive approach to eliminate memory safety vulnerabilities in C-based systems. This has even garnered attention from the US Department of Defense\footnote{\url{https://www.darpa.mil/news/2024/memory-safety-vulnerabilities}}, which has long-lived systems that disproportionately depend on C, supporting programs to translate C codebases to Rust (\href{https://www.darpa.mil/research/programs/translating-all-c-to-rust}{TRACTOR}). 
% 
Recent efforts like Syzygy~\citep{shetty2024syzygy}, C2SaferRust~\citep{nitin2025c2saferrust}, and AlphaTrans~\citep{alphatrans} have shown the potential for hybrid approaches combining LLMs with traditional program analysis techniques. However, some significant challenges remain, including ensuring correctness in large codebases while maintaining desirable attributes such as speed, reduced vulnerabilities, and idiomaticity.
% 
% \todo{expand a little bit, motivate real-world importance (DARPA, cobol to *), checked-c maybe}
% 
% \todo{
% \\	- Don't banks still use COBOL code for this reason? Too expensive to migrate (maybe because of legislative/bureaucratic overhead)?
% }


\textbf{Code Optimization:} Transforming programs to improve performance characteristics while maintaining functional correctness is a critical software task. Works have explored LLMs for optimizing simple standalone programs, such as PIE~\citep{pie} for C++ functions and AlphaDev~\citep{alphadev} to discover faster sorting algorithms at the assembly level. While these are considerably hard due to a large search space, optimizing real-world systems poses more significant challenges due to their scale, complexity, and competing objectives like speed and memory efficiency.
% 
% PIE
% LLM for compiler optimization \citep{llmcompiler}
% AlphaDev \citep{alphadev}

For instance, over two decades, changes to the Chrome web browser have been an exemplar of optimization affecting real-world code~\citep{chromium10years}. Their V8 engine achieved a 20x performance improvement through coordinated optimizations across multiple layers - from implementing concurrent garbage collection that reduced bloat by 100x to developing specialized compilers like TurboFan that improved performance by 5-10\%, to enabling background parsing and compilation that reduced compile time by 20\%. The demand for cross-layer and low-level code changes (e.g., writing a new JavaScript interpreter) and building tools to measure and test representative performance metrics are key challenges for such real-world impact with LLMs.

Additionally, the scope of code optimization can often be narrow, but the logical complexity can be extremely high. For example, optimizing kernel code at the PTX level for GPU-based AI model optimization can be very complex~\citep{deepep2025, ouyang2024kernelbench}. In many scenarios, high levels of autonomy may not be desirable, as tradeoffs can depend heavily on external factors such as hardware, and the best optimizations may ultimately affect readability.

% https://blog.chromium.org/2018/09/10-years-of-speed-in-chrome_11.html
% https://v8.dev/blog/ignition-interpreter


\subsection{Software Testing and Program Analysis}
Software inevitably has bugs of a wide variety in scope and logical complexity.
% 
For instance, minor correctness bugs might miss a few corner cases and require quick changes to patch~\citep{mosolygo2021rise}. However, complex concurrency bugs and security vulnerabilities can be deep in the call stack or hard to isolate and fix~\citep{concurrencylucene}.
% 
Software testing, both in development and production, is a practical approach to bug-free code. However, it faces challenges such as manual test creation and the scalability limits of traditional static analysis and testing tools. As LLMs continue to improve at code generation, researchers have explored using them to alleviate these challenges at different temporal horizons.

For short-horizon testing (e.g., unit tests~\citep{lemieux2023codamosa} and property-based tests~\citep{vikram2023can}), models have shown promise in generating targeted test cases \citep{mundler2025swt} and even hill-climb on metrics like code coverage~\citep{ryan2024code}. For simple, function-level testing, LLM-generated tests have been used to filter~\citep{chen2022codet, zhang2023algo} and debug programs \citep{chen2025revisit}. 
% 
% \todo{add more citations for llm testgen}
% 
Meta's ACH system~\citep{foster2025mutation}, a mutation-guided LLM-based test generator, was used to find undetected faults in Messenger and WhatsApp.
% 
In contrast, long-horizon testing approaches like fuzzing require sustained program state exploration, where recent work (Fuzz4All~\citep{xia2024fuzz4all}, KernelGPT~\citep{yang2023kernelgpt}, OSS-Fuzz~\citep{ossfuzzllm, ossfuzzllm-results}) shows LLMs can significantly improve effectiveness through better input generation and exploration strategies.

However, the most challenging security vulnerabilities and zero-day exploits, from memory corruption to privilege escalation, require deep program understanding that testing/fuzzing often miss.
% 
Works like IRIS~\citep{li2024llm}, InferROI~\citep{wang2023boosting}, and LATTE~\citep{liu2023harnessing} explored augmenting static analysis tools (e.g., CodeQL) with LLMs for taint and resource leak analyses.
% 
More recently, \citet{bigsleep} demonstrated the potential of using LLMs at a much bigger scale by finding a real SQLite vulnerability through exploratory variant analysis. That said, we are still scratching the surface - LLMs are still improving on real-world tasks that security researchers handle daily, like reasoning about complex program logic or autonomously conducting offensive security operations as measured by CyberSecEval~\citep{wan2024cyberseceval}.

% \todo{Some works about security:
% \\ SecCodePLT: A Unified Platform for Evaluating the Security of Code GenAI https://seccodeplt.github.io/}

Beyond long-running analysis to identify vulnerabilities, several traditional program analyses have struggled to scale in practice despite their theoretical promise~\cite{X}. For instance, inferring \textit{program invariants} (properties deemed to always be true at a program point) has been challenging with traditional symbolic methods such as Daikon \citep{ernst2007daikon, padon2016decidability} while being valuable for exposing bugs~\citep{hangal2002tracking} and aiding software evolution~\citep{ernst1999dynamically}. Similarly, type inference for dynamically typed languages suffers from coverage limitations of rule-based approaches and requires specialized tools like ShapeIt~\citep{zheng2024dynamic} for domain-specific challenges such as inferring symbolic tensor shapes.

AI models have addressed some of these constraints scalably across domains. For invariant detection, they enhance traditional approaches through structured representations~\citep{si2018learning}, LLM-based prompting~\citep{kamath2023finding} and re-ranking~\citep{chakraborty2023ranking}, and reinforcement learning~\citep{yu2023loop}. Similarly, works have used sequence-to-sequence models~\citep{wei2023typet5}, few-shot LLM approaches like TypeGen~\citep{peng2023generative}, and generate-then-rank methods like TIGER~\citep{wang2024tiger} for type inference.
% 
Consequently, we observe new benchmarks emerging in the space such as LIG-MM~\citep{liu2024towards} for loop-invariant detection. \todo{add more benchmarks in the space}


% \ks{Add a section on root cause analysis.}

\subsection{Code Maintenance}

Understanding the moving parts of a codebase is an important skill for being a stellar engineer. Code can often be difficult to understand due to many wrapper functions, error-handling boilerplate, deep call stacks, and sometimes even poor code cleanliness. We identify four practical code maintenance tasks that we find practical and important:

\textbf{Code Documentation and Summarization}: To ensure maintainability, readability, and ease of collaboration, code must be well documented. Good documentation needs to be written cleanly and crisply, describing both what the function does and how the function works. It must also anticipate and address any misunderstandings that a programmer might have, such as potential side effects or special cases. Documentation is generally a task that has a low logical complexity and does not require too much human intervention. Rather, humans often see documentation as a chore and neglect it, leading to code and documentation frequently being out of sync. This has led to the concept of ``self-documenting code'', code that clearly conveys its purpose. With AI, we believe that documentation can become a continuously updated artifact in sync with the code. 

% \ks{Another challenge in code documentation is that often the written code and documentation go out of sync.} 
Several works have used LLMs for code summarization invoking techniques like prompting \citep{sun2024source, su2024distilled, haldar2024analyzing, ahmed2024automatic}. RepoAgent \citep{luo2024repoagent} is a framework that analyzes global contextual relationships in source code to generate fine-grained documentation. \citet{shi2024natural} show that LMs are capable of generating good natural language outlines -- text descriptions alongside code to partition it into semantically coherent sections. One challenge is that the evaluation of this task is very tricky: the academic community currently lacks datasets and benchmarks that contain good documentation and the automatic evaluation metrics do not align well with human metrics \citep{diggs2024leveraging}. 

\textbf{Code Navigation}: It is often challenging to navigate and get accustomed to a mature codebase for new comers. Code navigation means finding where relevant functionality is implemented throughout a codebase. Doing this well requires understanding the high-level layout of where every functionality lies in the codebase and the low-level understanding of which helper functions are used to implement each functionality. 

\textbf{Pull Request (PR) Review}: To review pull requests, a person must be familiar enough with the codebase to judge how well newly introduced changes fit the style and structure of the existing code. The most essential requirement is that the PR implements the new feature correctly. In addition to this, the reviewer must also determine how to integrate the PR with the existing codebase more smoothly. This includes checking whether the repository's style conventions are satisfied, ensuring that the PR does not introduce any new bugs, verifying that program invariants and guarantees still hold, and inspecting whether tests are robust. All of these require a global understanding of the codebase. There have been a few works studying the ability of AI systems to automatically review PRs \citep{tufano2021towards, tufano2022using, li2022automating, li2024utilizing}. Recently, AutoCodeRover \citep{zhang2024autocoderover} combines LLMs with code search to automatically fix GitHub issues.


\textbf{Code Question Answering}: The holy grail of code understanding is the ability to answer arbitrarily complex questions about a codebase, which requires sophisticated code understanding and reasoning abilities. Models should not hallucinate or give incorrect information that skews a developer's mental model of the code. Beyond other tasks mentioned in this section, developers might commonly ask questions related to data flow (when and where data structures get mutated), code functionality (whether there are any side effects), performance characteristics (determining the runtime and memory complexity of a function), or error handling (whether certain corner cases are handled). 


\subsection{Program Repair}

% \citep{wang2023boosting}
% InferROI: LM to augment static analysis to detect resource leaks in Java programs

% IRIS \citep{li2024llm}: augment CodeQL with LM to improve taint analysis

% \textbf{Bug Detection}: Software inevitably will have bugs of a wide variety of flavors that vary in scope and logical complexity. Minor bugs might miss a few corner cases and require only a few lines of modification, e.g. \textit{simple, stupid bugs} \citep{mosolygo2021rise}. Complex bugs (e.g. concurrency bugs) might have very high logical complexity. Because bugs can often pass tests and syntax checks, even detecting them can be tricky and requires thorough testing both during development and production work. \todo{distribute across PL things / testing sub-sections and focus on only fixing parts in the sub-section}

Program repair is the task of repairing software bugs. The difficulty of fixing a bug depends on the scope and logical complexity\ks{What do you mean scope?  What is logical complexity?  Why not consider the logical complexity of the code?}. Function-level, low logical complexity bugs can easily fixed by feeding in the error information. Repository-level, high logical complexity bugs could require locating the bug, re-thinking the algorithm, and restructuring or refactoring the codebase. In addition, these changes must be general and should not break the functionality of other parts of the codebase.

Automated program repair has had a long history, with many benchmarks covering different scopes and languages. These include DroixBench \citep{tan2018repairing} for android apps; Defects4J \citep{just2014defects4j}, GitBug-Java \citep{gitbugjava}, and growingBugs \citep{GrowingBugsICSE21, GrowingBugsTSE2022, NaturalnessOfBugsFSE2022} for real-world Java; Bugsinpy \citep{widyasari2020bugsinpy} for Python; BugSwarm \citep{tomassi2019bugswarm} for multilingual; DebugBench \citep{hu2024leveraging}, LiveCodeBench \citep{jain2024livecodebenchholisticcontaminationfree}, and Codeflaws \citep{tan2017codeflaws} for LeetCode-style problems; and many more.

Historically, there have been many techniques for this task, including heuristic-based APR (using genetic programming to explore the search space of the correct patch), constraint-based APR (treating repair as a constraint-solving task), pattern-based APR (apply expert hand-crafted repair templates), and learning-based APR (using language models) \citep{zhang2024systematic}.
% \citep{jiang2023impact, xia2022practical, tufano2019empirical, haque2022fixeval, jin2023inferfix, gupta2017deepfix, berabi2021tfix}
More recently, with LLMs, there have been agent-based approaches such as FixAgent \citep{lee2024unified} using agents specializing in different aspects of debugging, and RepairAgent \citep{bouzenia2024repairagent} that invokes suitable tools. On the other hand, Agentless \citep{xia2024agentless} uses a three-phase process of localization, repair, and patch validation. \ks{All SWE agents should also be considered as repair agents?}

Finally, program repair has also been used as a tool to improve code generation, where error messages and incorrect test cases are fed back into the model to improve code generation \citep{madaan2023self, chen2024teaching, zhang2023self, olausson2024repair, zhong2024debug, tang2025code}. This is also known as self-repair/self-debugging. For a fantastic and much more comprehensive survey of automated program repair, we highly recommend the reader check out \href{https://program-repair.org/}{this website}\footnote{\url{https://program-repair.org/}}.


% multilingual \citep{liu2024mdeval}


% \citet{mosolygo2021rise} defines \textit{simple, stupid bugs} as ``minor inaccuracies that do not prevent the code from running or compiling, but change its behavior in a subtle way that may lead to unforeseen consequences''. 





\subsection{Scaffolding and Meta-Code} \label{subsec:scaffolding-metacode}

For a software system to work, the core logic must be written well, but that is not enough. Many infrastructural aspects must be in place to support the software. We group these into two main categories: \textit{scaffolding} and \textit{meta-code}. We define \textit{scaffolding} to be code that is important to make the system work but does not actually participate in the execution of its main logic. Examples of scaffolding include test harnesses, configuration files, CI/CD code, Makefiles, Dockerfiles, sandbox databases, and preprocessors. In contrast, we define \textit{meta-code} as a task outside of the code that must be done to get the software running the property. Examples of meta-code include setting up Google authentication, subscribing to APIs, managing file storage, and generating API tokens.\todo{It seems the first one is better called metacode and the second is better called scaffolding, no?}

\textit{Infrastructure-as-code and Security.} A particularly challenging case is generating \textit{Infrastructure-as-code} such as \texttt{Terraform}, where you specify the type of infrastructure specifications (\todo{examples}) as code and execute it to create the infrastructure. When generating such code, LLMs struggle with security configurations due to the complex interplay between service-level permissions (e.g., AWS resource access), resource-level permissions (e.g., specific allowed actions), and provider-specific security primitives like IAM roles, security groups, and network access controls.

\begin{tcolorbox}[colback=lightblue, boxrule=0pt, arc=10pt, outer arc=10pt]
\textit{Example. Distinguishing permission levels in cluster setup}: \cite{terrateam2024} show that on a task of bringing up a cluster, models fail to distinguish between ECS (Amazon Elastic Container Service) Task Execution Roles (for container operations) and Task Roles (for application-level permissions). This resulted in overly permissive policies where a single role was granted both image pull permissions and DynamoDB table access, violating principles of least privilege.
\end{tcolorbox}

\begin{tcolorbox}[colback=lightblue, boxrule=0pt, arc=10pt, outer arc=10pt]
\textit{Example. Configuration validation:} Ciri \citep{lian2024large} is a tool that uses LLMs for configuration validation on open-source projects including Django, PostgreSQL, and Redis. They find that while Ciri excels at detecting misconfigurations of syntax and range violations, it struggles to detect dependency and version violations and is limited to a narrow range of misconfigurations. They also find that LLMs are biased towards more popular configuration parameters, which may lead to hallucinations in out-of-domain scenarios. 
\end{tcolorbox}

% \todo{What other challenges here? Anything on test harnesses? examples from Devin blog?}

% \textit{Example. Playskript}: \href{https://code.playskript.com/}{Playskript} is a programming system that allows you to create interactive visualizations, presentations and games. AWS Lambda config, AWS virtual network + permissions, email service, google auth, databases, file storage. \todo{Ask Armando to flesh this out more}

\subsection{Formal Verification}

The task of formal verification involves generating checkable, mechanized proofs that can guarantee that a piece of code works as intended. Formal verification of software is necessary in mission-critical applications such as aircraft software, where it is crucial that code is correct with absolute certainty. In some cases, software bugs have led to costly disasters. In the maiden flight of the Ariane 5 rocket, a floating-point conversion error caused it to explode forty seconds after liftoff. Another case is with the computer-controlled radiation therapy machine Therac-25, where concurrency bugs led to six people being massively overdosed, leading to serious injury and deaths. 

In the formal methods literature, there are two major classes of formal verification: full functional verification (FFV) and property verification (PV). In FFV, the goal is to design a complete and precise formal specification that captures the desired behavior of the implementation. For example, \citet{zee2008full} developed a method for full functional verification of data structures like mutable lists, trees, graphs, and hash tables. The challenge in full functional verification generally lies in writing the specification in a way that captures all the required aspects that need to be shown.

While FFV provides a complete set of guarantees, it is usually sufficient to opt for PV, where a few key properties of a system are proven correct. Examples include: ensuring that two threads do not simultaneously enter a critical section of a program, verifying that a complex program will always terminate, proving the absence of security vulnerabilities like buffer overflows, and guaranteeing memory safety. Unlike in FFV, the specification is usually simpler to write, and the challenge is usually finding the right domain-specific tool and designing the right algorithms to write the proof.

As a field, formal software verification has seen a few great successes in the last few years: Astrée \citep{cousot2005astree} was able to completely verify that Airbus A340's primary flight-control software had no run-time errors, verifying 132,000 lines of C code. 
CompCert \citep{leroy2016compcert} is a formally verified optimizing C compiler that supports most of the ISO C 99 language. More recently, formal methods have been applied to verify a cryptographic server \citep{erbsen2024foundational} and an IoT lightbulb at both a hardware and software level \citep{erbsen2021integration}. At Amazon, formal methods been used to verify and protect cryptographic software \citep{goel2024leanprover}, cloud resources \citep{xu2024cloud}, and authorization \citep{disselkoen2025cedar}. For a survey containing a complete list of applications, we refer the reader to \citet{huang2023lessons}.

To make formal verification easier, there have been countless programming languages designed specifically for it. Some of the popular ones include TLA \citep{lamport1994introduction}, Coq \citep{Coq-refman}, Lean \citep{de2015lean}, Dafny \citep{leino2010dafny}, Isabelle \citep{nipkow2002isabelle}, and Verus \citep{lattuada2024verus}. 

A few works have explored using LLMs to write formal verification. Benchmarks like DafnyBench \citep{loughridge2024dafnybench} and miniCodeProps \citep{lohn2024minicodeprops} were designed to measure the ability of LLMs to write software proofs in Dafny and Lean, respectively. In Dafny, \citet{poesia2024dafny} use a combination of search and prompting to create a synthetic dataset of annotations greatly improving performance on DafnyBench. Clover \citep{sun2024clover} generates code alongside consistency checks (like Dafny annotations), \citet{li2025dafny} employ Dafny as an intermediate language to improve code generation, and \citet{misu2024towards} explore prompting and retrieval to generate Dafny. In Rust, Verus is a popular formal verification language, with AutoVerus \citep{yang2024autoverus} and AlphaVerus \citep{aggarwal2024alphaverus} generating verified specifications and proofs for Rust functions. Finally, \citet{silva2024leveraging}, Lean Copilot \citep{song2024towards}, and llmstep \citep{welleck2023llmstep} aim to provide IDE plugins to aid humans in writing code in Dafny and Lean.

% \textbf{Bottlenecks in LLMs for Software Verification}: While LLMs can sometimes verify simple function-level properties, they have not been able to scale beyond that due to a few fundamental challenges. First, languages used in formal verification such as Dafny and Verus are low-resource languages, and LLMs often produce code with incorrect syntax and semantics. Second, writing precise specifications that fully capture all the desired properties in a sound and complete manner can sometimes be difficult. Third, writing proofs can be tricky for LLMs as well. In verification languages, invariants often need to be written in a precise form, because tools like SMT solvers can rely on heuristics that are sensitive to the way that constraints are written.

